State Doob's martingale decomposition theorem

Let \( Y_t \in \mathcal{L}^1(Q) \) be adapted, then there exists a unique decomposition
\( Y_t = M_t - A_t \), where \( M_t \) is a \( Q \)-martingale and \( A \) is an adapted process
with \( A_0 = 0 \).
Proof:
\( A_t - A_{t-1} = -\mathbb{E}_Q[Y_t - Y_{t-1}|\mathcal{F}_{t-1}] \)
Then \( A \) is predictable and \( M_t = Y_t + A_t \) is a \( Q \)-martingale.

The continuous version (stated for completeness), the Doob-Meyer Decomposition theorem, is a little trickier:
A cadlag supermartingale \( Z \) with \( Z_0 = 0 \) and \( \mathcal{Z} = \{Z_\tau | \tau < \infty \} \),
where \( \tau \) is any stopping time, being uniformly integrable, has a unique decomposition 
 \( Z_t = M_t - A_t \) of a predictable \( A_t \) with \( A_0 = 0 \) and uniform integrable martingale \( M_t \).

Uniform integrability was the condition on a family of r.v. \( \mathcal{C} \), s.t. for all \( X \in \mathcal{C} \)
1. \( \exists M < \infty, \mathbb{E}[|X|] \leq M \)
2. \( \forall \varepsilon > 0, \exists \delta > 0, \forall A \in \mathcal{F}, Q[A] \leq \delta \implies \mathbb{E}[|X \mathbf{1}_A] \leq \varepsilon \).


Define a stopping time and state some facts 

A function \( τ : \Omega \mapsto \{0, 1, \dots, T \} \cup \{ \infty \} \) is called a stopping
time if \( \{ \tau = t \} \in \mathcal{F}_t \) for \( t \in \{0, 1, \dots, T \} \).
In particular \( \tau = c \) (constant) is a stopping time.
It is easy to show that \(\forall t, \{ \tau = t \} \in \mathcal{F}_t  \) is equivalent to the definition.
For \( 2 \) stopping times \( \tau, \sigma \) also \( \tau \wedge \sigma,\ \tau \vee \sigma, (\tau + \sigma) \wedge T \)
are stopping times and in discrete time also \( \tau\sigma \). However \( \tau - \sigma \) is not one.

Define the stopped martingale process \( Y^\tau_t(\omega) = Y_{t \wedge \tau}(\omega) \),
there holds the famous Doob's stopping theorem.

Let \( M \) be an adapted process s.t. \( M_t \in \mathcal{L}^1(Q) \) for each \( t \), then
the following is equivalent

1. \( M \) is a \( Q \)-martingale
2. For any stopping time \( \tau \) the stopped process \( M^\tau \) is a \( Q \)-martingale
3. \( \mathbb{E}_Q[M_{\tau \wedge T}] = M_0 \) for any stopping time.
(This generalizes to sub and super martingales)


What is the optimal stopping problem and the Snell envelope?

Let \( \mathcal{T} = \{\tau | \tau \text{ is a stopping time with } \tau \leq T\} \),
then the optimal stopping problem is
\( \max_{\tau \in \mathcal{T}} \mathbb{E}[H_\tau] \)
for some process \( H_t \) (usually some american claim).

Define \( U_T = H_T \) and \( U_t = H_t \wedge E[U_{t+1}|\mathcal{F}_t] \)
as the Snell envelope of \( H_t \).

Let \( H \) be an adapted process with \( \mathbb{E}_Q[\max_{t \leq T} |H_t| < \infty \). 
Then the Snell envelope \( U^Q \) of \( H \) with respect to \( Q \) is the smallest \( Q \)-supermartingale dominating \( H \):
If is another \( Q \)-supermartingale such that \( \tilde{U}_t \geq H_t \) \( Q \)-a.s., then \( \tilde{U}_t \geq U_t^Q \) \( Q \)-a.s. 
for all \( t \).

Let \( \tau^{u}_{\text{min}} := \text{min}\{t \geq u| U_t = H_t\}\) with \( \tau_{\text{min}} \iff u = 0 \), it holds
\( U_t = \mathbb{E}[H_{\tau^{u}_{\text{min}}}|\mathcal{F}_t] = \text{ ess sup }_{\tau \in \mathcal{T}_t} \mathbb{E}[H_\tau | \mathcal{F}_t ] \)
In particular, \( U_0 = \mathbb{E}[H_{\tau_{\text{min}}}] = \sup_{\tau \in \mathcal{T}} \mathbb{E}[H_\tau] \).

Thus we have a solution to our optimal stopping problem, actually this can be used to identify a solution: 
\( \tau \in \mathcal{T} \) is optimal iff. \( H_\tau = U_\tau \) \( P \)-a.s. and if \( U^\tau \) is a martingale.
Any optimal stopping time satisfies \( \tau \geq \tau_{\text{min}} \).

Now define the stopping time:
\( \tau_{\text{max}} = \inf \{ t \geq 0 | \mathbb{E}[U_{t+1} - U_t|\mathcal{F}_t] \neq 0 \} \wedge T = \inf \{t \geq 0 | A_{t+1} \neq 0 \} \wedge T \).
The stopping time \( \tau_{\text{max}} \) is the largest optimal stopping time. Moreover, a stopping time \( \tau \) is optimal if and only if \( P \) -a.s. 
\( \tau \leq \tau_{\text{max}} \) and \( U_\tau = H_\tau \).



What is the relation between optimal stopping times and hedging/execution strategies

Suppose an american type claim, then the optimal stopping problem is an optimal execution strategy for the buyer
of the claim. Conversely the snell envelope describes the optimal hedge.
More specifically,

Let \( H \) be the discounted american claim with snell envelope \( U \) with respect so some EMM \( Q \). Then there exists a \( d \)-dimensional predictable \( \chi \) 
s.t. 
\( U_u + \sum_{k=u+1}^t \chi_k \cdot (X_k - X_{k-1}) \geq H_t \) for \( t \geq u \). 
Any other dominating \( \tilde{U}_t \) for \( H_t \) satisfies \( \tilde{U}_t \geq U_t \) \( Q \)-a.s..
Note that this can be interpreted as the Doob martingale decomposition for \( H_t \) and if we set \( u = 0 \) and take expectations it follows 
that \( U_0 \) must be the arbitrage free price for \( H_t \) (it can be shown that \( \chi_k \) is self-financing).

\( H_\tau \leq U_0 + \sum_{k=1}^\tau \chi_k \cdot (X_k - X_{k-1}) \) for all \( \tau \in \mathcal{T} \) \( Q \).a.s
and equality holds iff. \( \tau \) is optimal with respect to \( Q \).

Let \( V_t = \mathbb{E}_Q[H_T|\mathcal{F}_t] \) then \( U_t \geq V_t \) for all \( t \) and if \( V_t \)
dominates \( H_t \) this is an equality.

Example for such domination are found if \( H_t \) is a submartingale.:
1. A convex \( f : \mathbb{R}^d \to [0, \infty) \) is applied 
to the discounted market \( X \). By Jensen's inequality
\( \mathbb{E}[f(X_{t+1})|\mathcal{F}_t] \geq f(E[X_{t+1}|\mathcal{F_t}] = f(X_t) \).
2. For an american call \( H_t = (X^1_t - \frac{K}{S^0_t})_{+} \) it holds for increasing \( S^0_t \) that
\( \mathbb{E}[H_{t+1}|\mathcal{F_t}] \geq H_t \), thus the Snell envelope coincides with \( V_t = \mathbb{E}_Q[H_T|\mathcal{F}_t] \)
so the arbitrage free price \( U_0 \) coincides with the European type. (The same doesn't apply to Puts or if dividends exist)


How's the Snell envelope constructed in the CRR model?

Let \( H_t = h_t(S_t) \) be the american claim contingent on \( S_t \), then
\( U_t = u(S_t) \) where
\( u_T(x) = h_T(x) \) and \( u_t(x) = h_t(x) \wedge (u_{t+1}(x\hat{b})q + u_{t+1}(x\hat{a})(1 - q) \)
then the space \( [0, T) \times [0, \infty) \) can be decomposed into the regions
\( \mathcal{R}_c = \{ (t, x) | u_t(x) > h_t(x) \} \quad \mathcal{R}_s = \{ (t, x) | u_t(x) = h_t(x) \} \)
and the minimal stopping time \( \tau_\text{min} \) can be described as first exit/entry time via 
\( \tau_\text{min} = \text{min} \{t \geq 0| (t, S_t) \nin \mathcal{R}_c \} = \text{min} \{t \geq 0 | (t, S_t) \in \mathcal{R}_s \} \).


Give the arbitrage bounds for american claims

Note that for a specific stopping time \( \tau \), \( H_\tau \) can be treated like an European claim
by noting that \( \tilde{H}_t = H_t \mathbf{1}_{\tau = t} \) is of european type with maturity \( T = \tau \).
The arbitrage free prices are then given by
\( \Pi(H_\tau) = \{ E[H_\tau|P \in \mathcal{P}, E[H_\tau] < \infty \} \).
Therefore it would make sense to define arbitrage prices \( \pi \) for discounted american claims like this
1. \( \pi \) is not too high, in the sense that there is \( \tau \in \mathcal{T} \) and \( \pi' \in \Pi(H_\tau) \) s.t. \( \pi \leq \pi' \)
2. The price \( \pi \) is not too low in the sense that there exists no \( \tau' \in \mathcal{T} \) such that \( \pi < \pi' \)
for all \( \pi' \in \Pi(H_{\tau'} \).
It can be shown more rigorously then for the set of equivalent martingale measures \( \mathcal{P} \)
\( \pi_{\text{inf}}(H) = \text{inf}_{P^\ast \in \mathcal{P}} \sup_{\tau \in \mathcal{T}} E_P[H_\tau] = \sup_{\tau \in \mathcal{T}} \inf_{P \in \mathcal{P}} E[H_\tau] \)
\( \pi_{\text{sup}}(H) = \text{sup}_{P^\ast \in \mathcal{P}} \sup_{\tau \in \mathcal{T}} E_P[H_\tau] = \sup_{\tau \in \mathcal{T}} \sup_{P \in \mathcal{P}} E[H_\tau] \)
Moreover, \( \Pi(H) \) either consists of one single point or does not contain its upper endpoint \( \pi_{\text{sup}}(H) \).

A claim is attainable if there's a \( \tau \in \mathcal{T} \) and self-financing trading strategy \( \chi \) s.t. it satisfies \( P \)-a.s.
\( V_t \geq H_t \) where \( V_t \) is the value process of \( \chi \).
In particular \( \Pi(H) \) is then a singleton and \( \pi_{\text{sup}}(H) \in \Pi(H) \) (in fact this is an equivalence).


